\documentclass{article}
\usepackage[UTF8]{ctex}
\usepackage{amsmath,amssymb,amsfonts}
\usepackage{enumitem}
\usepackage{geometry}
\geometry{a4paper, margin=1in}
% === 错题本排版优化（不参与题目标识） ===
\setlength{\parindent}{0pt}
\setlength{\parskip}{0.35em plus 0.10em minus 0.08em}
\setlength{\emergencystretch}{3em}
\setlist{itemsep=0.25em, topsep=0.40em, parsep=0pt}
\allowdisplaybreaks[2]
% === 错题本排版优化结束 ===
\begin{document}
\section*{电子科技大学2025年高等代数真题}
\section*{159．电子科技大学2025年高等代数试卷}
一．填空题．每题5分，共30分．\\

\subsection{题目1}
1．设2025阶方阵 \(A, B\) 的秩均为2020，则 \(A B\) 的秩 \(r(A B)\) 能取到的最小值为 \(\_\_\_\_\)\\
\(\_\_\_\_\)．．\\
\(\_\_\_\_\)

\subsection{题目2}
2．设 \(A, B, C, D, E, F\) 均为3阶方阵，且

\[
\left(\begin{array}{lll}
I_{3} & A & B \\
O & I_{3} & C \\
O & O & I_{3}
\end{array}\right)^{-1}=\left(\begin{array}{lll}
I_{3} & D & F \\
O & I_{3} & E \\
O & O & I_{3}
\end{array}\right)
\]

已知 \(A=\left(\begin{array}{lll}3 & 0 & 0 \\ 0 & 4 & 0 \\ 0 & 0 & 5\end{array}\right), B=\left(\begin{array}{lll}1 & 0 & 1 \\ 0 & 2 & 0 \\ 3 & 0 & 1\end{array}\right)\) ，且 \(C=A+B-I\) ，则行列式 \(\operatorname{det} F=\) \(\_\_\_\_\)．

\subsection{题目3}
3．设 \(A\) 为 4 阶方阵满足 \(A^{4}+3 A^{2}+2 I=3 A^{3}+3 A\) ，伴随阵 \((A-I)^{*},(A-2 I)^{*}\) 秩均为 0 ，则 \(A^{2}\) 的迹 \(\operatorname{tr}\left(A^{2}\right)=\)

\subsection{题目4}
4．多项式 \(x^{2025}+1\) 除以 \((x-1)^{2}\) 的余式 \(r(x)=\)

\subsection{题目5}
5．设正交矩阵 \(P=\left(\alpha_{1}, \alpha_{2}, \alpha_{3}, \alpha_{4}\right)\) ，若二次型 \(f\left(x_{1}, x_{2}, x_{3}, x_{4}\right)\) 在线性替换 \(X=P Y\) 下化为标准型\(3 y_{1}^{2}+2 y_{2}^{2}-y_{3}^{2}\) ，令 \(Q=\left(-\alpha_{3}, \alpha_{1}, \alpha_{2}\right)\) ，则 \(f\left(x_{1}, x_{2}, x_{3}\right)\) 在线性替换 \(X=Q Y\) 下化为 \(\_\_\_\_\)．

\subsection{题目6}
6．设复方阵 \(A\) 的初等因子组为 \(\lambda, \lambda^{2}, \lambda^{2},(\lambda-1)^{2}, \lambda+2,(\lambda+2)^{2}\) ，则 \(A\) 的最小多项式为 \(\_\_\_\_\)．
\end{document}
